\vspace{-0.5em}
\section{Discussion}\label{sec:discussion}
\vspace{-0.2em}
\subsection{Discussion of PaTAS Findings}

Beyond per-output assessment, PaTAS can estimate a static trust opinion of the network as \(T(\text{NN}) = PaTAS_{FF}((1,0,0))\), on the principle that a trustworthy model preserves trust. Under this definition the IPTA results of Experiment~\ref{exp:3} show clear degradation for adversarial samples (\cref{tab:mnistpoisipta}), a quantitative indicator of reduced reliability under local perturbation.

Although trust mass and accuracy are often correlated, our results reveal meaningful divergences. In the breast cancer task (\cref{tab:cancer_results}), a mixed trust profile \((0.25,0,0.75)\) on all features and labels yields a final trust mass of \(0.34\), higher than the fully uncertain-features/fully trusted-label case (\(0.32\)), yet the latter achieves the better test accuracy (96\% vs.\ 80\%). A similar divergence appears in the poisoned MNIST experiment (\cref{tab:mnistpoisipta}): clean digit \(3\) and clean digit \(6\) reach comparable accuracy, but IPTA assigns higher trust to \(3\) (\(0.878\)) than to \(6\) (\(0.866\)). Higher trust mass therefore does not imply superior predictive performance.

These divergences reflect that PaTAS evaluates the reliability of the inference path, not predictive accuracy alone. Digit \(6\) was a poisoned class during training, so the parameters involved in predicting it are less trusted than those for \(3\), even when the test samples themselves are clean. 

Interpretation of trust values, like accuracy, is application-dependent: \((0.4,0,0.6)\) may be inadequate in high-stakes settings yet sufficient elsewhere. PaTAS supports such reading with a single interpretable metric that complements accuracy, signaling fragility when accuracy appears high and stability when it slightly decreases.

\paragraph*{Obtaining input and label trust}\label{sec:acquisition}
A recurring practical question is where the input- and label-level trust opinions consumed by PaTAS come from. Three routes are available, in decreasing order of information. (i) \emph{Provenance-based assessment}: when data are drawn from sources of known reliability, the quantification models of \cref{sec:back} map the supporting and contradicting evidence for a source directly into a binomial opinion; this is how our companion work assesses dataset trustworthiness~\cite{10.1007/978-3-032-19096-3_17} and how feature-level opinions are assigned when part of an input, such as a region covered by an adversarial patch, is known to be unreliable. (ii) \emph{Estimation from data}: label agreement across annotators or models, and input-quality signals such as out-of-distribution or anomaly scores, can be converted to opinions through the same evidence mapping; the conformity-based construction of \cref{exp:4} instantiates this route and carries the detection results of \cref{tab:comparative}. (iii) \emph{Vacuous fallback}: when no such information exists, the opinion is initialized to \((0,0,1)\). Because a vacuous input yields a vacuous output (\cref{thm:infvac}), this fallback is conservative by construction: it produces uncertainty-dominated assessments rather than unwarranted trust, and PaTAS still contributes at inference time through parameter and path trust. Crucially, the per-query collection circumstances of a single input are usually easier to establish than the provenance of an entire training set, so input-level trust is often obtainable even when dataset-level trust is not.

A key factor is the threshold \(\epsilon\) separating positive from negative gradient evidence. As \cref{tab:cancer_results} shows, too small an \(\epsilon\) causes PaTAS to read gradients as large even when the model performs well, so that trusted labels reduce parameter trust, resembling the \(g \to +\infty\) case in Appendix~\ref{app:asymptotic}. Calibration therefore matters: \(\epsilon\) can be tuned per layer or tied to the learning rate, so that it decays alongside it and tracks the expected gradient scale throughout training.

\vspace{-0.5em}
\subsection{Threat Model and Security Scope}\label{sec:threat}
\vspace{-0.5em}

\paragraph*{Assets and goal}
The asset PaTAS protects is the \emph{reliability of individual model predictions at inference time}. A defender wishes to know, per query, whether an output can be trusted given the quality of the input, the reliability of the training data, and the parameters through which the input is propagated. PaTAS is a \emph{detection and monitoring} mechanism, not a hardening one; it does not alter the model or prevent an attack, but exposes predictions with weak supporting evidence.

\paragraph*{Adversary}
We consider a \emph{training-time data-poisoning adversary}~\cite{chen2017targetedbackdoorattacksdeep}. The adversary can inject or relabel a bounded fraction of training samples, including backdoor triggers such as fixed pixel patches, with the goal of inducing targeted misclassifications that remain hidden behind high nominal accuracy and confidence. The adversary does not control the model architecture, the training procedure, or the trust assessments assigned to known-provenance data. This is the threat evaluated in \cref{exp:3}.

\paragraph*{Defender knowledge}
PaTAS assumes white-box access to the model (parameters, activations, and gradients), as it runs alongside training and inference, and access to trust opinions on inputs and labels. When such opinions are unavailable, the defender initializes them to a vacuous opinion \((0,0,1)\), which yields conservative, uncertainty-dominated assessments rather than false assurance (\cref{sec:acquisition}).

\paragraph*{Out of scope}
Evasion attacks crafted at inference time, model-stealing, and membership-inference attacks are outside the scope of the present evaluation. PaTAS's trust and uncertainty signals may be applicable to them, for instance by monitoring the trust profile of query streams, but we make no such claim here and defer it to future work. We likewise prove consistency properties of the trust computation (\cref{sec:ptastheory}) rather than formal robustness guarantees against a bounded adversary, which remain an open direction.
